%% kapa: exam=1 round=28 year=2017 period=1 form=A subject=economics subject_ko=경제학원론 seq=30 number=70
%% source: Q-Net 감정평가사 기출문제 (cst00309) artlSeq=5199603 file=extracted/1008/5199603/2017년 제28회 감정평가사 1차시험 문제지/1교시 A형.pdf
%% conversion: kordoc 4.24.51 -> tools/gen1.py (auto); review=proofread-sonnet; flags=none
\begin{question}{70}
화폐수요함수는 $\frac{M^{d}}{P}=\frac{Y}{5i}$이다. 다음 중 옳은 것을 모두 고른 것은? (단, $\frac{M^{d}}{P}$는 실질화폐잔고, $i$는 명목이자율, $Y$는 실질생산량, $P$는 물가이다.)
\begin{framedtext}
ㄱ. 명목이자율이 일정하면, 실질생산량이 k\% 증가할 경우 실질화폐잔고도 k\% 증가한다. \par
ㄴ. 화폐유통속도는 $\frac{5i}{Y}$이다. \par
ㄷ. 명목이자율이 일정하면 화폐유통속도는 일정하다. \par
ㄹ. 실질생산량이 증가하면 화폐유통속도는 감소한다.
\end{framedtext}
\begin{choices}
  \item ㄱ, ㄴ
  \item ㄱ, ㄷ
  \item ㄴ, ㄷ
  \item ㄴ, ㄹ
  \item ㄷ, ㄹ
\end{choices}
\end{question}
\answer{2}
